If $m$ is even, each move changes an even number of orientations, so the parity of the number of upward cups is invariant. An initially odd number cannot reach zero; these pairs fail.

If $m$ is odd, any chosen pair of cups can be flipped while leaving the others unchanged: choose $m-1$ cups from the remaining $n-2$ (possible since $m<n$), and make the two moves consisting of these common cups together with either cup of the pair. The common cups flip twice. Therefore any even-sized set of cups can be flipped by pairing its members. To flip an odd-sized set, first make any legal odd-sized move, leaving an even-sized difference to correct. Hence every configuration is reachable.

For each $n$, precisely the odd values $1\le m<n$ work, numbering $\lfloor \dfrac{n}{2}\rfloor$. Thus the requested count is
\[
\boxed{\sum_{n=2}^{20}\left\lfloor\dfrac{n}{2}\right\rfloor
=2(1+2+\cdots+9)+10=100}.
\]
